% Generated by scripts/assemble.py; corrections belong in source/corrections.json.
\examhead{62}{重庆大学 2024 年数学分析试卷}{61}
% p062-b19; source PDF page 62
\par\Needspace{5\baselineskip}
\par\Needspace{6\baselineskip}
\blockanchor{p062-b19}
\majorquestion{一}计算数列极限\par
\nopagebreak[4]
% p062-b20; source PDF page 62
\blockanchor{p062-b20}
\displayformula{\lim _ {n \rightarrow \infty} \left(\frac {1}{\ln (1 + n) - \ln n} - n\right).}
% p062-b21; source PDF page 62
\par\Needspace{5\baselineskip}
\blockanchor{p062-b21}
\majorquestion{二}计算第一类曲面积分 $\iint_{S} (x^{2} + y^{2}) \, \mathrm{d}S$, 其中 $S$ 是立体 $\sqrt{x^2 + y^2} \leq z \leq 1$ 的边界.\par
% p062-b22; source PDF page 62
\par\Needspace{5\baselineskip}
\blockanchor{p062-b22}
\majorquestion{三}求级数 $\sum_{n=1}^{\infty} \frac{(-1)^n}{n(2n-1)}$ 的和.\par
% p062-b23; source PDF page 62
\par\Needspace{5\baselineskip}
\par\Needspace{6\baselineskip}
\blockanchor{p062-b23}
\majorquestion{四}计算第二类曲线积分\par
\nopagebreak[4]
% p062-b24; source PDF page 62
\blockanchor{p062-b24}
\displayformula{\int_ {L} \left(y ^ {2} - z ^ {2}\right) \mathrm{d} x + \left(2 z ^ {2} - x ^ {2}\right) \mathrm{d} y + \left(3 x ^ {2} - y ^ {2}\right) \mathrm{d} z.}
% p062-b25; source PDF page 62
\blockanchor{p062-b25}
其中 $L$ 是平面 $x + y + z = 2$ 与柱面 $|x| + |y| = 1$ 的交线，从 $z$ 轴正方向往下看去，$L$ 是顺时针方向.\par
% p062-b26; source PDF page 62
\par\Needspace{5\baselineskip}
\blockanchor{p062-b26}
\majorquestion{五}已知三元函数 $f(x, y, z)$ 在 $\mathbb{R}^3$ 上具有连续的二阶偏导数, 设 $\mathbb{R}^3$ 中光滑简单封闭曲面的全体为 $\Sigma$, 对于 $S \in \Sigma$, 用 $\pmb{n}$ 表示 $S$ 的外法线单位方向向量, $\frac{\partial f}{\partial \pmb{n}}$ 表示 $f$ 沿 $\pmb{n}$ 的方向导数.\par
% p063-b00; source PDF page 63
\blockanchor{p063-b00}
(1) 证明: $f$ 在 $\mathbb{R}^3$ 上满足\par
% p063-b01; source PDF page 63
\blockanchor{p063-b01}
$\frac{\partial^{2}f}{\partial x^{2}}+\frac{\partial^{2}f}{\partial y^{2}}+\frac{\partial^{2}f}{\partial z^{2}}=0$  (即 f 为  $R^{3}$  上的调和函数)\par
% p063-b02; source PDF page 63
\blockanchor{p063-b02}
当且仅当 $\iint_{S} \frac{\partial f}{\partial n} \, \mathrm{d}S = 0 (\forall S \in \Sigma)$.\par
% p063-b03; source PDF page 63
\par\Needspace{6\baselineskip}
\blockanchor{p063-b03}
(2) 设 $f(x, y, z)$ 为 $\mathbb{R}^3$ 上的调和函数, $S \in \Sigma$, 且 $S$ 围成的有界闭区域记作 $V$. 证明:\par
\nopagebreak[4]
% p063-b04; source PDF page 63
\blockanchor{p063-b04}
\displayformula{f (x _ {0}, y _ {0}, z _ {0}) = \frac {1}{4 \pi} \iint_ {S} \left(f \frac {\cos (\boldsymbol {r} , \boldsymbol {n})}{| \boldsymbol {r} | ^ {2}} + \frac {1}{| \boldsymbol {r} |} \frac {\partial f}{\partial \boldsymbol {n}}\right) d S.}
% p063-b05; source PDF page 63
\blockanchor{p063-b05}
其中 $(x_0, y_0, z_0)$ 为 $V$ 内部的一个定点，$\pmb{r} = (x - x_0, y - y_0, z - z_0)$，$(\pmb{r}, \pmb{n})$ 表示两向量的夹角，$|\pmb{r}|$ 表示向量 $\pmb{r}$ 的模长.\par
% p063-b06; source PDF page 63
\par\Needspace{6\baselineskip}
\blockanchor{p063-b06}
(3) 若函数 $f$ 和闭区间 $[0,1]$ 上的连续函数 $g$ 在单位球 $\Omega = \{(x,y,z) \mid x^2 + y^2 + z^2 \leq 1\}$ 上满足\par
\nopagebreak[4]
% p063-b07; source PDF page 63
\blockanchor{p063-b07}
\displayformula{\frac {\partial^ {2} f}{\partial x ^ {2}} + \frac {\partial^ {2} f}{\partial y ^ {2}} + \frac {\partial^ {2} f}{\partial z ^ {2}} = g (x ^ {2} + y ^ {2} + z ^ {2}).}
% p063-b08; source PDF page 63
\blockanchor{p063-b08}
证明： \displayformula{\iiint_{\Omega}\left(x\frac{\partial f}{\partial x} +y\frac{\partial f}{\partial y} +z\frac{\partial f}{\partial z}\right)\mathrm{d}x\mathrm{d}y\mathrm{d}z = \pi \int_0^1\sqrt{t} (t - 1)g(t)\mathrm{d}t.}\par
% p063-b09; source PDF page 63
\par\Needspace{5\baselineskip}
\blockanchor{p063-b09}
\majorquestion{六}判断下列四个命题的正误. 正确的, 请给出证明; 错误的, 需举出反例.\par
% p063-b10; source PDF page 63
\blockanchor{p063-b10}
(1) 若正项级数 $\sum_{n=1}^{\infty} a_n$ 收敛, 则 $\lim_{n \to \infty} \frac{a_{n+1}}{a_n} = 1$.\par
% p063-b11; source PDF page 63
\blockanchor{p063-b11}
(2) 若函数 $f(x)$ 在 $\mathbb{R}$ 上连续且有界, 则 $f(x)$ 在 $\mathbb{R}$ 上一致连续.\par
% p063-b12; source PDF page 63
\blockanchor{p063-b12}
(3) 若二元函数 $f(x, y)$ 在 (0,0) 点关于 $x$ 和 $y$ 都不存在偏导数，则 $f(x, y)$ 在点 (0,0) 沿 $x$ 轴的方向导数必不存在.\par
% p063-b13; source PDF page 63
\blockanchor{p063-b13}
(4) 若函数 $f(x)$ 在闭区间 $[a, b]$ 上有界且只存在有限个不连续点，则函数 $f(x)$ 在 $[a, b]$ 上Riemann可积.\par
% p063-b14; source PDF page 63
\par\Needspace{5\baselineskip}
\blockanchor{p063-b14}
\majorquestion{七}设参数 $\alpha > 0$, 讨论广义积分 $\int_{0}^{+\infty} \frac{e^{\sin x} \sin(2x)}{x^{\alpha}} \mathrm{d}x$ 的收敛性(包括条件收敛和绝对收敛).\par
% p063-b15; source PDF page 63
\par\Needspace{5\baselineskip}
\blockanchor{p063-b15}
\majorquestion{八}若函数 $f(x)$ 和 $g(x)$ 都在闭区间 $[a, b]$ 上连续, 在 $(a, b)$ 上二阶可导, 并且它们在 $(a, b)$ 内有相等的最大值, $f(a) = g(a)$, $f(b) = g(b)$, 求证: 存在 $\xi \in (a, b)$, 使得 $f''(\xi) = g''(\xi)$.\par
% p063-b16; source PDF page 63
\par\Needspace{5\baselineskip}
\blockanchor{p063-b16}
\majorquestion{九}函数列 $\{f_n(x)\}_{n=1}^{\infty}$ 在闭区间 $[a, b]$ 上逐点收敛于 $f(x)$.\par
% p063-b17; source PDF page 63
\blockanchor{p063-b17}
(1) 假设每个 $f_{n}(x)$ 都在 $[a, b]$ 上连续，并且 $\{f_{n}(x)\}_{n=1}^{\infty}$ 一致收敛，证明：$f(x)$ 在 $[a, b]$ 上连续.\par
% p063-b18; source PDF page 63
\blockanchor{p063-b18}
(2) 假设每个 $f_{n}(x)$ 都在 $[a, b]$ 上 Riemann 可积, 并且 $\{f_{n}(x)\}_{n=1}^{\infty}$ 一致收敛, 证明: $f(x)$ 在 $[a, b]$ 上 Riemann 可积.\par
% p063-b19; source PDF page 63
\blockanchor{p063-b19}
(3) 假设每个 $\{f_n(x)\}$ 都在 $[a, b]$ 上可微, 存在常数 $M > 0$, 对任意的 $x$ 和正整数 $n$, 都有 $|f_n'(x)| \leq M$. 证明: $\{f_n(x)\}_{n=1}^{\infty}$ 在 $[a, b]$ 上一致收敛.\par
